\subsection{群流形}

5.12.1. 设 G 是一个群。如果 G 上的流形结构与 G 的群结构相容，是指映射 \((x, y) \mapsto xy\) 从 \(G \times G\) 到 G 是态射。于是映射 \(x \mapsto x^{-1}\) 是从 G 到自身的态射。赋予 G 以其群结构和流形结构后，称其为群流形（若需精确说明，则称为“C＾\(C^{r}\)r 类群流形”），也称为李群。若 \(r = \omega\)，也称为解析群。若 K = R（分别为 C、\(Q_{p}\)），也称为实（分别为复、p-adic）李群。每个群流形都是纯的。群流形之间的同态（或简称同态）是一个从群流形到另一个群流形的映射，它既是群同态又是流形态射。

若 G 是群流形，则 G 的流形结构所对应的拓扑结构使 G 成为可度量化的完备拓扑群；当 K 局部紧且 G 有限维时，它是局部紧的。

5.12.2. 例：

\begin{enumerate}
\def\labelenumi{(\roman{enumi})}
\item
  若 V 是 Banach 空间，则 V 的典范流形结构与其交换群结构相容。
\item
  设 A 是带单位元的完备赋范 K-代数，A* 是 A 的可逆元素组成的群。它是 A 的一个开子空间，而由 K-向量空间 A 的典范流形结构在这个开集上诱导的流形结构与 A* 上的群结构相容。特别地，取 A 为 Banach 空间 E 的连续自同态代数 \(\mathcal{L}(E)\)；群 \(A^{*}\) 由拓扑向量空间 E 的自同构组成；把如此定义的群流形记为 \(\mathbf{GL}(E)\)。当 \(A = M_{n}(K)\) 时，我们在 \(\mathbf{GL}(n, K)\) 上得到群流形结构。
\item
  若 \(G_{1},\ldots,G_{n}\) 是群流形，则乘积群 \(G=G_{1}\times\cdots\times G_{n}\) 在赋予各 \(G_{i}\) 的流形结构的乘积流形结构后，是一个群流形。
\end{enumerate}

5.12.3. 设 G 是一个群流形，H 是一个拓扑群，且 \(f:H\to G\) 是满足第 5.8.1 号条件 (QR) 的连续同态。G 的流形结构在 \(f\) 下的逆像结构使 H 成为群流形。当 H 是 G 的一个子群且为子流形（分别为拟子流形）时尤其适用；这样的子群称为 G 的群子流形（分别为群拟子流形）；它必然是 G 中的闭集。

若 \(H_{i}(i=1,\ldots,n)\) 是 \(G_{i}\) 的群子流形，则 \(H_{1}\times\cdots\times H_{n}\) 是 \(G_{1}\times\cdots\times G_{n}\) 的群子流形。

5.12.4. 设 G 是一个群流形，H 是 G 的群子流形。等价关系 \(x^{-1}y \in H\) 是正规的，这使得可以在左陪集 xH 组成的空间 G/H 上赋予一个流形结构，称为 G 的流形结构的商。映射 \((g, x) \mapsto g.x\) 从 \(G \times (G/H)\) 到 G/H 是态射。对于由右陪集 Hx 组成的齐性空间 \(H\backslash G\)，也有类似结果。若 H 是正规子群，则 G/H 的流形结构与其群结构相容。

5.12.5. 设 G 是一个群流形，X 是一个流形。群 G 在流形 X 上的左作用是从 \((s, x) \mapsto sx\) 到 X 的态射 \(G \times X\)，满足

\[
\begin{array}{l l} s (t x) = (s t) x & \text { if } \quad s, t \in G, x \in X \\ e x = x & \text { if } \quad x \in X (e \text { being   the   identity   element   of } G). \end{array}
\]

也称群流形 G 在 X 上左作用；右作用律类似地定义。

令 \(x \in X\)，令 \(G_x\) 为 x 在 G 中的稳定子。假设映射 \(G\)\(g \mapsto g \cdot x\) 是子浸入（当 \(K\) 的特征为 0 且 X 有限维时，这一假设自动满足）。则 \(X\)\(G_x\) 是 G 的群子流形，由 \(G\) 通过取商得到的从 \(G/G_x\) 到 X 的映射是浸入。若进一步 x 的轨道 \(X\)\(g \mapsto g \cdot x\)\(G \cdot x\) 是局部闭的，且 G 的拓扑有可数基，则 \(x\)\(G\)\(G \cdot x\) 是 X 的子流形，而映射 \(X\)\(G/G_x \to G \cdot x\) 是流形同构。

